\documentclass[11pt]{article}

\usepackage[letterpaper,margin=0.85in]{geometry}
\usepackage[T1]{fontenc}
\usepackage{lmodern}
\usepackage{microtype}
\usepackage{amsmath,amssymb,mathtools,amsthm}
\usepackage{xcolor}
\usepackage[hidelinks]{hyperref}
\hypersetup{
  pdftitle={Consistency of Compressible Momentum Thickness},
  pdfauthor={Michael Gnadlinger},
  pdfsubject={Consistency proof for Schlichting's compressible momentum-thickness definition},
  pdfkeywords={compressible boundary layer, momentum thickness, Schlichting, typographical error}
}
\urlstyle{same}

\definecolor{HypBlue}{HTML}{17324D}
\definecolor{HypSteel}{HTML}{456C86}
\definecolor{HypLightBlue}{HTML}{EAF2F7}
\definecolor{HypLightGreen}{HTML}{EAF6EF}
\definecolor{HypLightRed}{HTML}{FBEDEC}
\definecolor{HypGray}{HTML}{F5F7F9}
\definecolor{HypRule}{HTML}{D8E0E6}

\setlength{\parindent}{0pt}
\setlength{\parskip}{0.65em}
\setlength{\emergencystretch}{2em}

\newcommand{\note}[2]{%
  \vspace{0.3em}%
  \noindent\colorbox{#1}{%
    \begin{minipage}{0.965\linewidth}
      #2
    \end{minipage}}%
  \vspace{0.3em}%
}

\newcommand{\eqbox}[2]{%
  \vspace{0.1em}%
  \noindent\fcolorbox{HypRule}{#1}{%
    \begin{minipage}{0.955\linewidth}
      #2
    \end{minipage}}%
  \vspace{0.1em}%
}

\newcommand{\dd}{\mathop{}\!\mathrm{d}}
\newcommand{\rhoe}{\rho_e}
\newcommand{\ue}{u_e}
\newcommand{\taUw}{\tau_w}
\newcommand{\theC}{\theta_{\mathrm c}}
\newcommand{\theS}{\theta_{\mathrm S}}
\newtheorem{proposition}{Proposition}

\title{\vspace{-3em}\textcolor{HypBlue}{Consistency of Compressible Momentum Thickness Formula in Schlichting's Bounadry-Layer Theory}\\[1.2em]
\large\textit{A short proof concerning Schlichting's Eqs.~(10.92) and (10.96)}\vspace{-1em}}
\author{Michael Gnadlinger}
\date{August 7, 2026}

\begin{document}
\maketitle
\vspace{-1em}
\hrule
\vspace{-1em}

\section{Summary}
\label{sec:summary}

This report concerns H.~Schlichting and K.~Gersten's \emph{Boundary-Layer
Theory}, 9th ed. (Springer, 2017), specifically Eqs.~(10.92) and
(10.96)~\cite{Schlichting2017}. The momentum thickness printed as Eq.~(10.96),
denoted by $\theS$ below, is
\[
 \theS=\int_0^h\frac{\rho u}{\rhoe\ue}
 \left(1-\frac{\rho u}{\rhoe\ue}\right)\dd y.
\]
When read literally, this expression is not generally compatible with the
momentum-integral equation printed as Eq.~(10.92) for a variable-density
profile. With the standard zero-defect normalization, compatibility instead
requires
\[
 \delta_2=\theC=
 \int_0^h\frac{\rho u}{\rhoe\ue}
 \left(1-\frac{u}{\ue}\right)\dd y.
\]
Importantly, the same corrected definition appears in similar foundational
books by White and Majdalani, Eq.~(7-75), and Anderson,
Eq.~(6.137)~\cite{White2022,Anderson2006};
Section~\ref{sec:cross-check} discusses this independent cross-check.

\note{HypLightBlue}{%
\begin{itemize}
    \item \textbf{Inconsistency localized:} the second density factor in the
    deficit term of the printed Eq.~(10.96) changes the transported momentum
    flux from $\rho u^2$ to $\rho^2u^2/\rhoe$.
    \item \textbf{Compatible definition:} the given momentum-integral equation Eq.~(10.92) requires the conventional
    velocity-deficit form of the momentum thickness, shown above, once the common zero-defect normalization
    is imposed.
\end{itemize}
}

\section{Consistency proof}
\label{sec:proof}

\subsection*{Scope}
Schlichting's momentum-integral equation, Eq. (10.92), is accepted as given,
including the assumptions under which it was derived. The purpose is only to
determine which profile integral can represent the momentum thickness in that
equation. A fixed upper limit $h$ is taken in uniform edge flow; it may equally
be replaced by the boundary-layer edge or by infinity.

\eqbox{HypLightGreen}{%
\textbf{Result.}
With the standard zero-defect normalization, the momentum thickness compatible
with Schlichting's Eq. (10.92) is
\begin{equation}
 \displaystyle
 \delta_2=\theC=
 \int_0^h\frac{\rho u}{\rhoe\ue}
 \left(1-\frac{u}{\ue}\right)\dd y.
 \label{eq:conventional}
\end{equation}
The expression printed as Schlichting's Eq. (10.96), if read literally, is not
generally compatible with Eq. (10.92) unless an additional density-defect term
is retained.
}

\begin{proof}
Schlichting's momentum-integral equation Eq. (10.92) is
\begin{equation*}
 \delta_2'
 +\frac{\delta_2}{\ue}\ue'
 \left(2+\frac{\delta_1}{\delta_2}-M_e^2\right)
 =\frac{\taUw}{\rhoe\ue^2}.
 \tag{S10.92}
\end{equation*}
The origin of the $-M_e^2$ term follows from the edge-flow relations. For an
inviscid edge flow, streamwise momentum and isentropic thermodynamics give
\[
 p_e'=-\rhoe\ue\ue',
 \qquad
 p_e'=a_e^2\rhoe',
\]
where $a_e$ is the edge speed of sound. Equating these expressions and using
$M_e=\ue/a_e$ yields
\begin{equation}
 \frac{\rhoe'}{\rhoe}
 =-\frac{\ue\ue'}{a_e^2}
 =-M_e^2\frac{\ue'}{\ue}.
 \label{eq:edge}
\end{equation}
Schlichting and Gersten state this exact edge relation explicitly as their
Eq.~(10.100) and note that it follows from the momentum equation together with
the condition $T_0=\mathrm{const}$ at the outer edge of the boundary
layer~\cite{Schlichting2017}. Thus, Eq.~\eqref{eq:edge} is not an additional
relation introduced here, but the outer-flow relation used in Schlichting's
derivation of Eq.~(10.92).
Accordingly, the $-M_e^2$ coefficient in Eq. (S10.92) accounts for the streamwise
variation of the edge density.

Multiplying Eq. (S10.92) by $\rhoe\ue^2$ and collecting the $\delta_1$ term
separately gives
\begin{equation*}
 \taUw
 =\rhoe\ue^2\left[
   \delta_2'+(2-M_e^2)\frac{\delta_2}{\ue}\ue'
  \right]
 +\rhoe\ue\delta_1\ue'.
\end{equation*}
The bracketed terms form a product derivative. Indeed, the product rule gives
\begin{align*}
 \frac{\dd}{\dd x}(\rhoe\ue^2\delta_2)
 &=\rhoe\ue^2\left[
   \delta_2'
   +\left(\frac{\rhoe'}{\rhoe}
          +2\frac{\ue'}{\ue}\right)\delta_2
  \right]\\
 &=\rhoe\ue^2\left[
   \delta_2'+(2-M_e^2)\frac{\delta_2}{\ue}\ue'
  \right],
\end{align*}
where the second line uses Eq.~\eqref{eq:edge}. Substitution into the preceding
equation therefore gives

\eqbox{HypGray}{%
\begin{equation}
 \displaystyle
 \taUw=\frac{\dd}{\dd x}(\rhoe\ue^2\delta_2)
       +\rhoe\ue\delta_1\ue'.
 \label{eq:Sconservative}
\end{equation}
}

It remains to identify the first term physically. Define
\[
 I_m=\int_0^h\rho u\dd y,
 \qquad
 I_p=\int_0^h\rho u^2\dd y.
\]
Integration of steady mass and streamwise momentum conservation through the
boundary layer gives
\begin{equation}
 I_m'+\rhoe v_e=0,
 \qquad
 I_p'+\rhoe\ue v_e=-h p_e'-\taUw.
 \label{eq:integrated}
\end{equation}
Using the inviscid edge relation $p_e'=-\rhoe\ue\ue'$ to eliminate $p_e'$, and
using the first of Eqs.~\eqref{eq:integrated} to eliminate $v_e$, yields
\begin{equation}
 \taUw=\frac{\dd}{\dd x}(\ue I_m-I_p)
       +\ue'(\rhoe\ue h-I_m).
 \label{eq:raw}
\end{equation}
Now introduce Schlichting's displacement thickness, Eq. (10.95),
\begin{equation}
 \delta_1=\int_0^h
 \left(1-\frac{\rho u}{\rhoe\ue}\right)\dd y,
 \label{eq:delta1}
\end{equation}
and note that
\begin{equation}
 \ue I_m-I_p
 =\int_0^h\rho u(\ue-u)\dd y\equiv J.
 \label{eq:J}
\end{equation}
Equation~\eqref{eq:raw} therefore becomes the physical momentum-deficit balance

\eqbox{HypGray}{%
\begin{equation}
 \displaystyle
 \taUw=J'+\rhoe\ue\delta_1\ue'.
 \label{eq:physical}
\end{equation}
}

Equations~\eqref{eq:Sconservative} and \eqref{eq:physical} contain identical
wall-stress and displacement-thickness terms. Hence
\begin{equation}
 \frac{\dd}{\dd x}(\rhoe\ue^2\delta_2-J)=0.
 \label{eq:comparison}
\end{equation}
The standard normalization that both quantities vanish with the momentum
deficit removes the possible $x$-independent constant. Thus
\[
 \rhoe\ue^2\delta_2=J
 =\int_0^h\rho u(\ue-u)\dd y,
\]
which, after division by $\rhoe\ue^2$, proves Eq.~\eqref{eq:conventional}.

For comparison, a literal reading of Schlichting's Eq. (10.96) defines
\begin{equation*}
 \theS=\int_0^h\frac{\rho u}{\rhoe\ue}
 \left(1-\frac{\rho u}{\rhoe\ue}\right)\dd y.
 \tag{S10.96}
\end{equation*}
Multiplying the two candidate definitions by $\rhoe\ue^2$ gives
\begin{align}
 \rhoe\ue^2\theC
 &=\int_0^h(\rho u\ue-\rho u^2)\dd y,
 \label{eq:correctflux}\\
 \rhoe\ue^2\theS
 &=\int_0^h\left(\rho u\ue-\frac{\rho^2u^2}{\rhoe}\right)\dd y.
 \label{eq:wrongflux}
\end{align}
The momentum equation transports $\rho u^2$, as present in
Eq.~\eqref{eq:correctflux}, not $\rho^2u^2/\rhoe$ as present in
Eq.~\eqref{eq:wrongflux}. Therefore Eq. (S10.96) is not the momentum-deficit
coefficient in Eq.~\eqref{eq:physical} for a general variable-density profile.
\end{proof}

\section{Consequence and checks}
\label{sec:consequence}

\subsection{Exact consequence of using Eq.~(S10.96)}
Let
\begin{equation}
 \Delta\theta\equiv\theS-\theC
 =\int_0^h\frac{\rho u}{\rhoe\ue}\frac{u}{\ue}
 \left(1-\frac{\rho}{\rhoe}\right)\dd y.
 \label{eq:difference}
\end{equation}
Since $\theC=\theS-\Delta\theta$, the correct balance becomes
\begin{equation}
 \taUw=\frac{\dd}{\dd x}(\rhoe\ue^2\theS)
 +\rhoe\ue\delta_1\ue'
 -\frac{\dd}{\dd x}(\rhoe\ue^2\Delta\theta).
 \label{eq:corrected}
\end{equation}
Equivalently, Eq. (10.92) written with $\theS$ requires the residual

\eqbox{HypLightRed}{%
\begin{equation}
 \displaystyle
 \mathcal R_S=\frac{1}{\rhoe\ue^2}
 \frac{\dd}{\dd x}(\rhoe\ue^2\Delta\theta)
 =\Delta\theta'+(2-M_e^2)\frac{\ue'}{\ue}\Delta\theta.
 \label{eq:residual}
\end{equation}
}

Constant density is a sufficient special case for $\Delta\theta=0$; it is not
an assumption made in the proof.

\subsection{Assumptions used}
Equation (S10.92) and its original assumptions are accepted. The short
derivation of Eq.~\eqref{eq:physical} uses the same standard ingredients:
steady two-dimensional thin-boundary-layer conservation, a stationary
impermeable wall, a uniform edge state, no streamwise body force, and the
inviscid edge momentum relation. Equation~\eqref{eq:edge}, already encoded in
Eq. (S10.92), is the additional isentropic edge relation needed for its
Mach-number form. No constant-density, zero-pressure-gradient, or thermal-wall
assumption is introduced.

\subsection{Cross-check with Schlichting}
\label{sec:cross-check}
Equation~\eqref{eq:delta1} is exactly Schlichting's Eq. (10.95), while
Eq. (S10.92) is his Eq. (10.92). His Eqs. (10.97) and (10.98) also retain the
expected pattern of a mass-flux weight multiplied by a deficit in the
transported quantity. Independent standard references print
Eq.~\eqref{eq:conventional}: White and Majdalani, Eq. (7-75), and Anderson,
Eq. (6.137)~\cite{White2022,Anderson2006}. These checks localize the
inconsistency to the second density factor in the printed Eq. (10.96). The
argument proves incompatibility of the two printed equations; whether the
cause is a typographical error is an editorial question.

\begin{thebibliography}{9}
\small

\bibitem{Schlichting2017}
H. Schlichting and K. Gersten,
\emph{Boundary-Layer Theory}, 9th ed., Springer, 2017,
Eqs. (10.92)--(10.100).

\bibitem{White2022}
F. M. White and J. Majdalani,
\emph{Viscous Fluid Flow}, 4th ed., McGraw Hill, 2022,
Sec. 7-4.2, Eq. (7-75), p. 418.

\bibitem{Anderson2006}
J. D. Anderson, Jr.,
\emph{Hypersonic and High-Temperature Gas Dynamics}, 2nd ed., AIAA, 2006,
Sec. 6.7.7, Eq. (6.137), p. 333.

\end{thebibliography}

\end{document}
